\chapter{State of the Art and Literature Review}
\label{chap:state-of-the-art}

This chapter shows what is already known and where your project fits. It should compare sources, methods, tools, standards, or systems rather than simply list them.

\section{Theoretical Background}
Introduce the concepts required to understand the project. Define the technical vocabulary, models, standards, or scientific principles that will be used later.

Avoid explaining every basic term from the field. Focus on background that directly supports the problem, method, or evaluation.

\section{Existing Methods and Technologies}
Describe the main approaches currently used to solve similar problems. For each approach, explain its principle, strengths, limits, and conditions of use.

\placeholderfigure{literature-map.pdf}{Literature map placeholder}{Insert a diagram that groups the main families of methods, tools, or research directions.}

\section{Related Systems}
Compare representative systems, tools, frameworks, products, or academic prototypes. Use consistent comparison criteria so the reader can see why your choices are reasonable.

\begin{table}[H]
    \centering
    \caption{Comparison matrix template}
    \label{tab:comparison-matrix-template}
    \small
    \begin{tabular}{L{0.22\textwidth}L{0.22\textwidth}L{0.22\textwidth}L{0.24\textwidth}}
        \toprule
        \textbf{Approach} & \textbf{Strengths} & \textbf{Limitations} & \textbf{Relevance to this project} \\
        \midrule
        Approach A & Replace with evidence-based strengths. & Replace with observed or cited limits. & Explain the lesson retained. \\
        \tabrowsep
        Approach B & Replace with evidence-based strengths. & Replace with observed or cited limits. & Explain the lesson retained. \\
        \bottomrule
    \end{tabular}
\end{table}

Avoid unsupported rankings such as ``best'' or ``most powerful'' unless the criteria and evidence are explicit.

\section{Identified Gaps}
End the chapter by stating the gap that your project addresses. The gap may be technical, methodological, organizational, experimental, or usability-related.

Avoid forcing a gap that your work does not actually address. A modest, well-supported gap is stronger than an exaggerated claim.

\section{Conclusion}
Summarize the lessons taken from the literature and explain how they shape the requirements or methodology of the next chapters.